\section*{Executive summary}

Who bears an AI labour shock determines what it costs. This study builds the first open-source, full tax-benefit microsimulation of AI displacement incidence for the United Kingdom, and the first anywhere to treat \emph{who} loses work --- not only how many --- as a first-class scenario axis. Occupation-level AI exposure is attached to workers in the Family Resources Survey, displacement and wage shocks are simulated across 66 shock-size scenarios and four incidence families, and the resulting household incomes are passed through the full UK tax and benefit system using PolicyEngine UK. The incidence gradient in each family is an assumption, stated plainly as such; the contribution of the paper is to show how the tax-benefit system mediates each assumed incidence into fiscal, poverty and inequality outcomes.

Holding the aggregate shock fixed at its central calibration --- 7 per cent of employees displaced and a 2.6 per cent wage gain for those who remain --- the annual Exchequer cost ranges from \pounds 14.2 billion when displacement is uniform across workers to \pounds 24.5 billion under an expertise-compression incidence in which experienced, higher-paid workers bear the shock. The exposure-proportional family, our central calibration, costs \pounds 18.4 billion (Monte Carlo mean \pounds 18.0 billion, SD \pounds 1.5 billion across displacement draws); a junior-concentrated incidence costs \pounds 19.4 billion. Poverty moves less across families than the fiscal cost does: before-housing-costs poverty rises by 1.8 to 1.9 percentage points in three of the four families and by 2.1 points under expertise compression --- the family that is simultaneously the most expensive for the Exchequer and the most poverty-raising. Top-loaded incidence is fiscally expensive; flatter or junior-concentrated incidence is cheaper for the Exchequer but leaves poverty roughly as high.

One result is robust to every assumption we vary: inequality rises. The Gini coefficient of equivalised household disposable income increases in all 66 shock-size cells (by 0.06 to 1.55 percentage points) and in all four incidence families at the central shock (by 0.92 to 1.14 points). This holds even though the market shock itself is top-loaded in the exposure-proportional family: the share of workers at risk of displacement rises monotonically from 0.5 per cent in the bottom decile of household income to 5.1 per cent in the top. A shock whose direct incidence falls mainly on better-off households still raises measured inequality once losses propagate through households and the tax-benefit system. Who bears the shock decides whether the UK faces primarily a fiscal problem or a poverty problem; it does not decide whether inequality rises.

In the central scenario, 1.61 million employees are displaced, before-housing-costs poverty rises by 1.87 percentage points, and the Gini rises by 1.10 points from a baseline of 0.309. Lost income tax and National Insurance receipts, together with higher benefit spending, outweigh the additional revenue from wage gains. Wage gains for workers who remain are nearly flat across the distribution (2.5 to 2.7 per cent by decile), so they do little to offset the distributional effects of displacement.

Two policy implications follow. First, the composition of support should track the incidence that materialises: a top-loaded shock strains the revenue base and calls for fiscal buffers, while a junior- or uniform-incidence shock calls for poverty-focused transfers and retraining that reaches workers before displacement. Second, because inequality rises under every incidence we consider, distributional monitoring cannot wait for the incidence question to be settled empirically --- the direction of the inequality effect is already clear even where its size is not.
